\documentclass[11pt,a4paper]{article}
\usepackage[utf8]{inputenc}
\usepackage[T1]{fontenc}
\usepackage[margin=2cm, top=2.4cm]{geometry}
\usepackage{xcolor}
\usepackage{enumitem}
\usepackage{titlesec}
\usepackage{fancyhdr}
\usepackage{hyperref}
\usepackage{booktabs}
\usepackage{tabularx}
\usepackage{verbatim}
\usepackage{amsmath,amssymb}

% Course color palette
\definecolor{mlblue}{HTML}{0066CC}
\definecolor{mlpurple}{HTML}{3333B2}
\definecolor{mllavender}{HTML}{ADADE0}
\definecolor{mllavender2}{HTML}{C1C1E8}
\definecolor{mllavender3}{HTML}{CCCCEB}
\definecolor{mllavender4}{HTML}{D6D6EF}
\definecolor{mlorange}{HTML}{FF7F0E}
\definecolor{mlgreen}{HTML}{2CA02C}
\definecolor{mlred}{HTML}{D62728}
\definecolor{mlgray}{HTML}{7F7F7F}
\definecolor{lightgray}{HTML}{F0F0F0}
\definecolor{midgray}{HTML}{B4B4B4}

% Section heading colors
\titleformat{\section}{\large\bfseries\color{mlpurple}}{}{0em}{}[\vspace{-4pt}\textcolor{mllavender}{\rule{\linewidth}{1pt}}\vspace{2pt}]
\titleformat{\subsection}{\normalsize\bfseries\color{mlblue}}{}{0em}{\textbullet~}
\titlespacing*{\section}{0pt}{8pt}{4pt}
\titlespacing*{\subsection}{0pt}{5pt}{2pt}

% Header/footer
\pagestyle{fancy}
\fancyhf{}
\setlength{\headheight}{22pt}
\addtolength{\topmargin}{-10pt}
\renewcommand{\headrulewidth}{0.4pt}
\renewcommand{\headrule}{\hbox to\headwidth{\color{mllavender}\leaders\hrule height \headrulewidth\hfill}}
\fancyhead[L]{\small\color{mlpurple}\textbf{Layer 2 Scaling Solutions} \textcolor{mlgray}{|} Lesson 10 \textcolor{mlgray}{|} Pre-Class Discovery Handout}
\fancyhead[R]{\small\color{mlgray}Prof.\ Dr.\ Joerg Osterrieder}
\fancyfoot[C]{\small\color{mlgray}\thepage}

% Compact list spacing
\setlist[itemize]{nosep, topsep=2pt, partopsep=0pt, leftmargin=1.4em}
\setlist[enumerate]{nosep, topsep=2pt, partopsep=0pt, leftmargin=1.6em}

% Activity box helper
\newcommand{\activitybox}[3]{%
  \noindent\colorbox{mllavender4}{\parbox{\dimexpr\linewidth-2\fboxsep}{%
    \textbf{\color{mlpurple}#1}\hfill\textcolor{mlgray}{\small #2}\\[2pt]#3%
  }}\par\vspace{4pt}%
}

% Thin ruled cell for fill-in tables
\newcommand{\fillcell}{\rule{2.2cm}{0.3pt}}

\begin{document}

\begin{center}
{\LARGE\color{mlpurple}\textbf{Layer 2 Scaling Solutions}}\\[4pt]
{\large\color{mlblue}Pre-Class Discovery Handout}\\[2pt]
{\small\color{mlgray}Lesson 10 $\cdot$ Complete before class $\cdot$ 25--30 minutes}
\end{center}
\vspace{4pt}

% -- Activity 1 ----------------------------------------------------------------
\activitybox{Activity 1: Explore L2 Solutions}{10 min}{%
Visit \textbf{L2Beat.com}, the canonical aggregator of Layer 2 data, and explore the current L2 ecosystem. Answer:
\begin{enumerate}
\item What is the total TVL across all Layer 2 solutions? How does it compare to Ethereum L1 DeFi TVL?
\item List the top 3 L2s by TVL. For each note: name, TVL, technology type (Optimistic / ZK / Other), and TPS.
\item Pick one L2 and examine its ``Risk Analysis'' page on L2Beat. What are the main trust assumptions? Is the sequencer centralized?
\item Compare the average transaction fee on Ethereum L1 vs.\ your chosen L2. What is the approximate cost ratio?
\end{enumerate}
\textbf{Bonus:} Find an L2 with a ``Stage~2'' designation on L2Beat. What does Stage~2 mean for decentralization and user security?
\vspace{4pt}
{\footnotesize\renewcommand{\arraystretch}{1.35}%
\setlength{\tabcolsep}{5pt}%
\begin{tabular}{@{}lllll@{}}
\toprule
\textbf{L2 Name} & \textbf{TVL (\$)} & \textbf{Type} & \textbf{Avg Fee (\$)} & \textbf{Stage} \\
\midrule
\fillcell & \fillcell & \fillcell & \fillcell & \fillcell \\
\fillcell & \fillcell & \fillcell & \fillcell & \fillcell \\
\fillcell & \fillcell & \fillcell & \fillcell & \fillcell \\
\bottomrule
\end{tabular}}
\vspace{2pt}
}

\vspace{2pt}

% -- Activity 2 ----------------------------------------------------------------
\activitybox{Activity 2: Rollup Comparison}{10 min}{%
Investigate the two dominant rollup architectures and compare their key properties. Answer:
\begin{enumerate}
\item How does an Optimistic Rollup handle an invalid transaction? How long does the challenge period typically last, and what happens during that window?
\item How does a ZK Rollup prevent invalid transactions from being accepted? What is the role of the prover, and why is no challenge period needed?
\item Which rollup type achieves faster finality on L1? Why does the other type require a delay?
\item Which type is currently cheaper to operate, and why? What is the main computational bottleneck for the other type?
\end{enumerate}
\vspace{4pt}
{\footnotesize\renewcommand{\arraystretch}{1.35}%
\setlength{\tabcolsep}{5pt}%
\begin{tabular}{@{}p{3.4cm}p{4.6cm}p{4.6cm}@{}}
\toprule
\textbf{Dimension} & \textbf{Optimistic Rollup} & \textbf{ZK Rollup} \\
\midrule
Proof Type           & \fillcell & \fillcell \\
Finality Time        & \fillcell & \fillcell \\
Computation Cost     & \fillcell & \fillcell \\
EVM Compatibility    & \fillcell & \fillcell \\
Major Examples       & \fillcell & \fillcell \\
Trust Assumption     & \fillcell & \fillcell \\
\bottomrule
\end{tabular}}
\vspace{2pt}
}

\vspace{2pt}

% -- Activity 3 ----------------------------------------------------------------
\activitybox{Activity 3: Bridge Security Analysis}{5 min}{%
Research major bridge hacks to understand the security risks of cross-chain asset transfers. Answer:
\begin{enumerate}
\item Look up the \textbf{Ronin Bridge hack} (March 2022). How much was stolen? What was the root cause (validator key compromise, smart contract bug, or other)?
\item Look up the \textbf{Wormhole hack} (February 2022). What specific vulnerability was exploited in the smart contract?
\item Estimate the total value lost across the top 5 bridge hacks in history. What percentage of current total L2 TVL does this represent?
\item What is the difference between a \textit{trusted bridge} and a \textit{trustless bridge}? Which offers stronger security guarantees and why?
\end{enumerate}
\vspace{4pt}
\noindent
\textbf{Total Bridge Losses:}~\$\fillcell\quad
\textbf{Largest Single Hack:}~\$\fillcell\quad
\textbf{Root Cause Pattern:}~\fillcell
\vspace{2pt}
}

\vspace{2pt}

% -- Activity 4 ----------------------------------------------------------------
\activitybox{Activity 4: L2 Fee Analysis}{5 min}{%
Compare transaction costs across Ethereum L1 and major L2 networks. Answer:
\begin{enumerate}
\item Go to \textbf{l2fees.info}. Compare the cost of a simple ETH transfer on: Ethereum L1, Arbitrum, Optimism, zkSync Era, and Base. Record your findings in the table below.
\item How did \textbf{EIP-4844} (Dencun upgrade, March 2024) change L2 fees? What are ``blobs'' and why do they reduce L2 costs by 10--100$\times$?
\item If a user makes 10 transactions per day, calculate the monthly cost savings of using an L2 vs.\ Ethereum L1.
\item How does the L2 sequencer earn revenue? How does the total user fee get split between L1 data posting costs and L2 execution costs?
\end{enumerate}
\vspace{4pt}
{\footnotesize\renewcommand{\arraystretch}{1.35}%
\setlength{\tabcolsep}{5pt}%
\begin{tabular}{@{}llll@{}}
\toprule
\textbf{Chain} & \textbf{ETH Transfer (\$)} & \textbf{Token Swap (\$)} & \textbf{Monthly Cost (10 tx/day)} \\
\midrule
Ethereum L1  & \fillcell & \fillcell & \fillcell \\
Arbitrum     & \fillcell & \fillcell & \fillcell \\
Optimism     & \fillcell & \fillcell & \fillcell \\
zkSync Era   & \fillcell & \fillcell & \fillcell \\
Base         & \fillcell & \fillcell & \fillcell \\
\bottomrule
\end{tabular}}
\vspace{2pt}
}

\vspace{4pt}

% -- Glossary ------------------------------------------------------------------
\section{Key Terms}

\renewcommand{\arraystretch}{1.25}
\noindent\begin{tabular}{@{}p{3.6cm}p{11.0cm}@{}}
\toprule
\textbf{Term} & \textbf{Definition} \\
\midrule
\textbf{Layer 2 (L2)}               & A scaling solution built on top of a Layer 1 blockchain (like Ethereum) that processes transactions off-chain while inheriting the security of the underlying L1 through proofs or fraud challenges. Reduces congestion and fees without changing the base layer. \\
\textbf{Rollup}                     & A Layer 2 scaling technique that executes transactions off-chain, then posts compressed transaction data back to L1. Anyone can verify correctness using this data. The two main variants are optimistic rollups and zero-knowledge rollups. \\
\textbf{Optimistic Rollup}          & A rollup that assumes all transactions are valid by default and only checks them if a verifier submits a fraud proof during the challenge period (typically 7 days). Examples include Arbitrum, Optimism, and Base. \\
\textbf{Zero-Knowledge Rollup}      & A rollup that generates a cryptographic validity proof (ZK-SNARK or ZK-STARK) for every batch of transactions, providing instant mathematical certainty of correctness. No challenge period is needed. Examples: zkSync Era, StarkNet, Polygon zkEVM. \\
\textbf{Fraud Proof}                & A mechanism in optimistic rollups where any verifier can challenge an invalid state transition by submitting evidence to the L1 contract, triggering re-execution. If fraud is confirmed, the sequencer is penalized and the state is corrected. \\
\textbf{Validity Proof}             & A cryptographic proof (ZK-SNARK or ZK-STARK) submitted by a ZK rollup prover that mathematically guarantees a batch of transactions was executed correctly. Accepted by the L1 contract without trusting the prover. \\
\textbf{Sequencer}                  & The entity responsible for ordering transactions, executing them off-chain, and submitting batches to L1. Most current L2 designs use a single centralized sequencer, which is a key trust and liveness assumption. \\
\textbf{State Channel}              & A Layer 2 technique where participants open a channel on-chain, conduct many off-chain transactions between themselves, then settle the final state on-chain. Best suited for repeated interactions between fixed parties (e.g., the Lightning Network for Bitcoin). \\
\textbf{Sidechain}                  & An independent blockchain with its own consensus mechanism, connected to a main chain via a bridge. Unlike rollups, sidechains do \emph{not} inherit L1 security---they rely on their own validator set. Example: Polygon PoS. \\
\textbf{Bridge}                     & A protocol enabling asset transfers between different blockchains or between L1 and L2. Bridges lock assets on the source chain and mint equivalent tokens on the destination chain. They are a major attack surface: over \$1B has been stolen from bridges. \\
\textbf{Data Availability (DA)}     & The guarantee that transaction data is published and accessible so anyone can independently verify the L2 state. Without data availability, users cannot reconstruct the L2 state or challenge invalid transitions. A critical security property for rollups. \\
\textbf{Blob (EIP-4844)}            & A new transaction type introduced in Ethereum's Dencun upgrade (March 2024) providing cheap, temporary data storage for rollups. Blobs are pruned after $\approx$18 days and reduced L2 costs by 10--100$\times$ compared to posting data as calldata. \\
\textbf{ZK-SNARK}                   & Zero-Knowledge Succinct Non-Interactive Argument of Knowledge. A proof system producing small, fast-to-verify proofs. Requires a trusted setup ceremony to generate public parameters. Used by zkSync Era and Polygon zkEVM. \\
\textbf{ZK-STARK}                   & Zero-Knowledge Scalable Transparent Argument of Knowledge. A proof system requiring no trusted setup, making it fully transparent and quantum-resistant. Produces larger proofs than SNARKs but is verifiable without any ceremony. Used by StarkNet and StarkEx. \\
\bottomrule
\end{tabular}

\vspace{6pt}
\noindent\textcolor{mlgray}{\small\textit{Prepared by Prof.\ Dr.\ Joerg Osterrieder \,\textbullet\, Layer 2 Scaling Solutions --- Lesson 10 \,\textbullet\, Pre-Class Discovery Handout}}

\end{document}
